\subsection{标量本质可积函数的积分的性质}

命题 5. --- 设 \(\mu\) 为 T 上的有界正测度，S 为承载 \(\mu\) 的 \(\mu\) -可测集（Ch. V, §5, No.~7, 定义 4），f 为取值于 F 中的标量 \(\mu\) -可积 (*) 函数。设 D 为 \(\mathbf{f}(\mathrm{S})\) 在空间 \(\mathrm{F}'^*\)（赋以拓扑 \(\sigma(\mathrm{F}'^*, \mathrm{F}')\)）中的闭凸包。则 \(\int \mathbf{f} d\mu \in \mu(\mathrm{T})\mathrm{D}\)。

由于 D 是包含 \(\mathbf{f}(\mathrm{S})\) 的诸闭半空间之交（TVS, II, §5, No.~3, 命题 4 的推论 1），只需证明：关系 \(\langle\mathbf{f}(t),\mathbf{z}^{\prime}\rangle\leqslant a\)（对一切 \(t\in S\)；其中 \(z^{\prime}\in F^{\prime},a\in R\)）蕴含 \(\langle z^{\prime},\int f d\mu\rangle\leqslant a\cdot\mu(T)\)；而由于 \(\int f d\mu=\int_{S}f d\mu\)，这由 Ch. IV, §4, No.~2, 定理 1 的推论 1 即得。

推论. --- 设 \(\mu\) 为 T 上的有界正测度，S 为承载 \(\mu\) 的 \(\mu\) -可测集，f 为 T 到 F 中的映射，它标量 \(\mu\) -可测且使 \(\mathbf{f}(\mathbf{S})\) 含于 F 的一个弱紧凸子集 A。则 f 是标量 \(\mu\) -可积的，且 \(\int f d\mu \in \mu(T)A \subset F\)。

对每个 \(z^{\prime} \in F^{\prime}\)，\(\langle z^{\prime}, f \rangle\) 在 S 中 \(\mu\) -可测且有界，故可积，这就证明 f 标量可积。此外，由于 A 在 \(F_{\sigma}\) 中紧，它在 \(F^{\prime*}\) 中闭，且 \(\mathbf{f}(S)\) 在 \(F^{\prime*}\) 中的闭凸包含于 A，推论得证。

命题 6. --- 设 f 为取值于 F 中的标量本质 μ-可积函数，使 \(\int f d\mu \in F\)。对 F 上每个下半连续半范数 q，5

\[
q \left(\int \mathbf {f} d \mu\right) \leqslant \int^ {\bullet} (q \circ \mathbf {f}) d \mu .
\]

设 D 为 \(z \in F\) 中使 \(q(\mathbf{z}) \leqslant 1\) 的点所成之集；D 闭、凸且含 0，故 \(D = D^{\circ\circ}\)（TVS, II, §6, No.~3, 定理 1 的推论 3）。因此只需证明：对每个 \(z' \in D^{\circ}\)，有 \(\left|\left\langle z', \int f d\mu \right\rangle\right| \leqslant \int^{\bullet}(q \circ f) d\mu\)；而这由 \(|\langle z', f(t) \rangle| \leqslant q(f(t))\)（对每个 \(t \in T\)）这一事实立即得出。

注意数值函数 \(q \circ f\) 不一定 \(\mu\) -可测（习题 12）。

命题 7. --- 设 f 为 T 到 F 中的标量本质 \(\mu\) -可积映射，使对 T 的每个紧子集 K，\(\mathbf{f}(\mathbf{K})\) 含于 F 的一个弱紧均衡凸子集。则 \(\int f d\mu\) 属于 F 的双对偶 \(F''\)。

对每个紧子集 \(\mathrm{K}\)（\(\mathrm{T}\) 的），

\[
\int \mathbf {f} \varphi_ {\mathrm{K}} d \mu = \int (\mathbf {f} \varphi_ {\mathrm{K}}) d (\varphi_ {\mathrm{K}} \cdot \mu);
\]

命题 5 的推论可以应用于有界测度 \(\varphi_{\mathrm{K}} \cdot \mu\) 与函数 \(\mathbf{f}\varphi_{\mathrm{K}}\)，因此 \(\int \mathbf{f}\varphi_{\mathrm{K}} d\mu \in \mathrm{F}\)。对每个 \(\mathbf{z}' \in \mathrm{F}'\)，\(\langle \mathbf{z}', \mathbf{f} \rangle\) 是本质 \(\mu\) -可积的，于是（Ch. V, §1, No.~3, 命题 10）

\[
\int \langle {\bf z} ^ {\prime}, {\bf f} \rangle d \mu = \lim _ {\mathrm{K}} \int \langle {\bf z} ^ {\prime}, {\bf f} \rangle \varphi_ {\mathrm{K}} d \mu ,
\]

极限是关于 T 的紧子集的递增有向集取的。由此得出，关于这个集，\(\int f\varphi_{K}d\mu\) 收敛于 \(\int f d\mu\)（依拓扑 \(\sigma(F'^{*},F')\)）。现在，

\[
\left| \left\langle \mathbf {z} ^ {\prime}, \int \mathbf {f} \varphi_ {\mathrm{K}} d \mu \right\rangle \right| = \left| \int \langle \mathbf {z} ^ {\prime}, \mathbf {f} \rangle \varphi_ {\mathrm{K}} d \mu \right| \leqslant \int | \langle \mathbf {z} ^ {\prime}, \mathbf {f} \rangle | d \mu ,
\]

这证明元素 \(\int f\varphi_{K}d\mu\) 所成之集是 \(F_{\sigma}\) 的有界子集，从而也是 F 的有界子集（TVS, IV, §1, No.~1, 命题 1）。因此命题 7 是下述引理的推论：

引理 1. --- \(\mathbf{F}^{\prime*}\) 中（关于拓扑 \(\sigma(\mathbf{F}^{\prime*},\mathbf{F}^{\prime})\)）\(\mathbf{F}\) 的每个有界子集的闭包都含于双对偶 \(\mathbf{F}''\)。

因为 F 的有界子集含于 \(F''\) 中的这样一个极集——F 的强对偶 \(F'\) 中 0 的某邻域的极集——故在 \(F''\) 中关于 \(\sigma(\mathrm{F}'',\mathrm{F}')\) 相对紧（TVS, III, §3, No.~5, 命题 7 与 No.~4, 命题 4 的推论 2）；由于 \(\sigma(\mathrm{F}'',\mathrm{F}')\) 由 \(\sigma(\mathrm{F}'*,\mathrm{F}')\) 诱导，引理得证。

推论. --- 假设 F 半自反，并设 f 为 T 到 F 中的标量本质 \(\mu\) -可积映射，使对 T 的每个紧子集 K，\(\mathbf{f}(\mathbf{K})\) 有界。则 \(\int f d\mu\) 属于 F。

因为 F 的每个有界子集都相对弱紧（TVS, IV, §2, No.~2, 定理 1），且 \(F = F''\)。

命题 8. --- 设 \(\mu\) 为 T 上的有界正测度，S 为承载 \(\mu\) 的 \(\mu\) -可测集，f 为 T 到 F 中的 \(\mu\) -可测映射，使 f(S) 含于 F 的一个完备、有界、均衡凸子集 B。则 f 是标量 \(\mu\) -可积的，且 \(\int f d\mu \in \mu(T)B \subset F\)。

由于 S 是 \(\mu\) -可积的，存在 S 的一个分割，由一个 \(\mu\) -可忽略集 N 与一列紧子集 \((\mathrm{K}_{n})\) 组成，使 f 在每个 \(K_{n}\) 上的限制连续（Ch. IV, §4, No.~6, 定理 4 的推论 3 与 §5, No.~1, 定义 1）；因此 \(\mathbf{f}(\mathrm{K}_n)\) 是 F 的紧子集。闭均衡凸包 \(\mathbf{B}_n\)（\(\mathbf{f}(\mathrm{K}_n)\) 的）于是是预紧的（TVS, II, §4, No.~1, 命题 3）且含于 F 的完备子集 B，故它是紧的，从而更是弱紧的。因此（命题 5 的推论）\(\mathbf{f}\varphi_{\mathrm{K}_n}\) 是标量 \(\mu\) -可积的，且

\[
\mathbf {z} _ {n} = \int \mathbf {f} \varphi_ {\mathrm{K} _ {n}} d \mu \in \mu (\mathrm{K} _ {n}) \mathrm{B} _ {n} \subset \mu (\mathrm{K} _ {n}) \mathrm{B}.
\]

对 F 上每个连续半范数 p，由此得出

\[
p (\mathbf {z} _ {n}) \leqslant \mu (\mathrm{K} _ {n}) \cdot \sup _ {\mathbf {x} \in B} p (\mathbf {x});
\]

由于 B 有界，且以 \(\mu(\mathrm{K}_{n})\) 为通项的级数收敛且有和 \(\mu(\mathrm{T})\)，可见以 \(s_{n}=z_{1}+z_{2}+\cdots+z_{n}\) 为通项的序列是 F 的完备子集 \(\mu(\mathrm{T})\mathrm{B}\) 中的 Cauchy 序列。因此这个序列收敛于 \(\mu(\mathrm{T})\mathrm{B}\) 的一个元素 s；由于可设 \(\mathbf{f}(t)=0\) 在 T-S 上成立，对每个函数 \(\langle z', f\rangle\)（\(z'\in F'\)）应用 Lebesgue 定理就证明 \(s=\int f d\mu\)。
% semantic-fragments: legacy-input-seam
\relax{}
